\section{Preliminary width and tilt estimates}\label{app:preliminary-estimates}

The estimates here use only the tents at $\pi/8,\pi/4,3\pi/8$.
They apply to a full niche and to any larger partial barrier. In a
partial-turn comparison a horizontal reflection also reflects the
actual barrier; it need not produce another partial-turn objective.

Throughout this appendix put
\[
 r=\sqrt2,\quad k=r-1,\quad c=\cos(\pi/8),\quad s=\sin(\pi/8).
\]
Useful exact identities and bounds are
\begin{equation}\label{eq:preliminary-constants}
 \begin{gathered}
 k^2=1-2k,\quad s/c=k,\quad k+1/k=2r,\quad
 1/s+1/c=4c,\quad 1/s-1/c=4s,\\
 \frac25<k<\frac{29}{70},\quad k<\frac{169}{408},\quad s>\frac38.
 \end{gathered}
\end{equation}
The rational square-root comparisons follow by squaring positive
sides; for example $577^2-2\cdot408^2=1$ proves the last upper
bound for $k$.

\subsection{Short caps, including a positive width gap}

\begin{lemma}[Short three-tent bound]\label{lem:preliminary-short}
Let $U$ have height one and projection $[-a,a]$, with
$0<a\le1001/1000$. Suppose its roof $A$ is affine on
$J=[-a/2,a/2]$, with $A(-a/2)=1-h$, $A(a/2)=1$, $0\le h<1$.
If a nonnegative barrier $B$ dominates the three actual tents, then
\[
 \int_{[-a,a]\setminus J}A-\int_JB<\frac{41}{50}.
\]
\end{lemma}
\begin{proof}
Raise the left exterior roof by $h$, put the middle roof at height
one, and retain the right roof. The resulting roof $\widehat A$ is
concave: its left slopes are nonnegative and its right slopes are
nonpositive. It defines a height-one cap $\widehat U$, and its
exterior area exceeds the original by $ah/2$.

The first-quarter supports are unchanged, since the old point
$(a/2,1)$ dominates every raised left point in those directions.
If $\widehat G(t)=h_{\widehat U}(t+\pi/2)$, then
\[
 h_U(t+\pi/2)=
 \max\{\widehat G(t)-h\cos t,\ \cos t-(a/2)\sin t\}.
\]
The wall of the second term is nonpositive on $J$. Thus the
positive two-wall tent of $U$ is obtained from that of $\widehat U$
by lowering its companion wall by $h$.

Put
\[
 u=rh_{\widehat U}(\pi/4)-1,\quad
 v=rh_{\widehat U}(3\pi/4)-1,\quad M_0=(u+v)/2.
\]
The two middle points give $a/2\le u,v\le a$. The lifted cap is
below $A_0(x)=\min(1,1+u-x,1+v+x)$. Write
$D=\int_{[-a,a]\setminus J}(A_0-\widehat A)\ge0$.
The diagonal tent is $(\min(u-k-x,v-h-k+x))_+$.
If the proposed score is at least $41/50$, its exterior area
already gives $a\ge41/50$ and $h\le362/1001<a/2$, so the tent's
apex lies in $J$. Exact integration then gives
\begin{equation}\label{eq:short-relaxation}
 \begin{split}
 P\le F_h(a,u,v)-D,\qquad
 F_h={}&a-\frac{(a-u)^2+(a-v)^2}{2}-\frac{ah}{2}\\
 &-(M_0-h/2-k)_+^2\\
 &+\frac{(u-k-a/2)_+^2+(v-h-k-a/2)_+^2}{2}.
 \end{split}
\end{equation}
The preliminary restrictions also give $a-h>k$. For fixed $a,h$
in this competitive range, $F_h$ is concave on $[a/2,a]^2$. In the
positive-apex region its Hessian has diagonal entries $-3/2$
plus the corresponding clipping indicator, and off-diagonal
entries $-1/2$. In the other region the clipping terms vanish
and the Hessian is $-\operatorname{Id}$. The first derivatives
agree across the interfaces. Its unclipped interior critical point
is $u=v=(a+k+h/2)/2$. Therefore, with
$B_0(a)=a-(a-k)^2/2$,
\begin{equation}\label{eq:short-diagonal-maximum}
 F_h\le B_0(a)-kh/2-h^2/8.
\end{equation}
Here $B_0$ increases up to $a=1+k$. Since $B_0(9/10)<41/50$,
the hypothetical high score forces $a>9/10$. Also
\[
 B_0(1001/1000)=\frac{2001}{1000}k-\frac1{2000000}.
\]
At $h=1/23$, using $k<29/70$, the right side of
\cref{eq:short-diagonal-maximum} is less than
\[
 \frac{6071014497}{7406000000}<\frac{41}{50}.
\]
Consequently
\begin{equation}\label{eq:short-stability}
 a>9/10,\qquad h<1/23,\qquad D<9/1000.
\end{equation}

Set $t_0=323/1000$. Maximizing \cref{eq:short-relaxation}
under $u\ge a/2+t_0+h/4$ gives the constrained critical point
\[
 u=a/2+t_0+h/4,\qquad
 v=a/2+(2k-t_0)/3+h/4.
\]
Both clips vanish there. The $v$ derivative is zero and the $u$
derivative is $(2k-4t_0)/3<0$, so concavity proves this is the
constrained maximum. Its value is at most
\[
 B_0(1001/1000)-\frac23(t_0-k/2)^2
 =\frac{7649}{3000}k-\frac{1417319}{6000000}<\frac{41}{50}.
\]
The same calculation with the two coordinates exchanged gives
\[
 u-a/2<t_0+h/4,\qquad v-a/2<t_0+h/4.
\]
The KKT points lie in the square because $a>9/10$ and $h<1/23$.
These inequalities remove both clips in the actual diagonal tent.
The bounds $a>9/10$ and $h<1/23$ also make its apex positive.

Put $m=M_0-h/4$. In this regime
\[
 F_h=a-(a-m)^2-(m-k)^2-\frac{(u-v)^2}{4}-kh/2-h^2/8.
\]
The completion $a-(a-m)^2\le m+1/4$ excludes $m\le11/20$.
For $11/20<m\le16/25$, the symmetric part increases with both
$a\le1001/1000$ and $m\le16/25$, and its upper value is
\[
 \frac{82}{25}k-\frac{538921}{1000000}
 <\frac{5739553}{7000000}<\frac{41}{50}.
\]
Thus $m>16/25$.

We now pay for two genuine extra tents. Define the lifted support
deficits
\[
 \delta_L=c+sv-h_{\widehat U}(5\pi/8),\qquad
 \delta_R=c+su-h_{\widehat U}(3\pi/8).
\]
They are nonnegative and at most $s(v-a/2),s(u-a/2)$.
A support cut at $(-v,1)$ has horizontal legs
$\delta_L/s$ and $\delta_L/(c-s)$. Its whole charged triangle
has area $K\delta_L^2$, where $K=1/k$.
If the latter leg exceeded $a-v$, the cut with
$\delta_0=(c-s)(a-v)$ would already give
$D\ge(a-v)^2/r$. But
\[
 a-v>a/2-t_0-h/4>23/200,
 \qquad (23/200)^2/r>3703/400000>9/1000,
\]
contradicting \cref{eq:short-stability}. The right cut is identical.
They lie in separate wings, so
\begin{equation}\label{eq:short-cut-payment}
 D\ge K(\delta_L^2+\delta_R^2).
\end{equation}

Put
\[
 \begin{gathered}
 \eta=2/25,\quad \Gamma=Kc^2=\frac{4+3r}{4},\quad
 \eta_L=\eta-(1-k)h,\\
 E_L=(\eta_L-\delta_L/c)_+,\quad E_R=(\eta-\delta_R/c)_+.
 \end{gathered}
\]
The identities in \cref{eq:preliminary-constants} imply
\[
 1-1/c+k^2>\eta,\qquad 1/c+K\eta<32/25.
\]
At the left zero $x_0=k-v+h$ of the diagonal tent, the actual
$\pi/8$ companion wall has height at least $E_L$.
A point attaining $u$ has $x+y=1+u$ and $y\le1$, and lies on
the unchanged right wing. Its first-wall lower bound is
$1-1/s+(u-x)/k$. This dominates the additional ascending wall
up to $x_0+E_L/(1-k)$, because
\[
 2M_0-h>32/25-h/2>1/c+KE_L.
\]
That endpoint is before the diagonal apex: use
$k+\eta/(1-k)<99/175$ and $m>16/25$.
In the reflected coordinate, an attaining point for $v$ lowered
by $h$ supplies the other companion wall. The corresponding
condition is $2M_0-kh>32/25>1/c+KE_R$.
This gives two disjoint additional triangles, with full areas
$\Gamma E_L^2,\Gamma E_R^2$. Only their exterior tips can leave
$J$. Each omitted horizontal length is at most
$L_0+h/4$, where
\[
 L_0=t_0-k+\eta/k=\frac{483-920k}{1000}.
\]
The two omitted areas sum to at most $k(L_0+h/4)^2$. Hence
\[
 \int_J B\ge N_{45}+\Gamma(E_L^2+E_R^2)-k(L_0+h/4)^2,
\]
where $N_{45}$ is the diagonal-tent integral used above.
This inequality remains valid when one of the peaks is zero.

For $\eta_i,\delta\ge0$, completing a square gives
\[
 K\delta^2+\Gamma(\eta_i-\delta/c)_+^2\ge\Gamma\eta_i^2/2.
\]
Combine this with \cref{eq:short-cut-payment} to obtain
\[
 P\le B_0(1001/1000)-kh/2-h^2/8
       -\frac\Gamma2\bigl[(\eta-(1-k)h)^2+\eta^2\bigr]
       +k(L_0+h/4)^2.
\]
Its coefficient of $h$ is $(1403/2000)k-19/50<0$, and its
quadratic coefficient is $-1/8-\Gamma(1-k)^2/2+k/16<0$.
Its largest value for $h\ge0$ is consequently at zero, and equals
\[
 \frac{8238929}{1000000}k-\frac{5185441}{2000000}
 <\frac{334549037}{408000000}<\frac{41}{50}.
\]
The final gap is $10963/408000000>0$, contradicting the high
score and proving the lemma.
\end{proof}

\subsection{Wide caps and the endpoint equations}

\begin{lemma}[Wide three-tent bound]\label{lem:preliminary-wide}
Let $8/5\le a<3$, $0\le h<1/2$, and let $U,A,J,B$ satisfy the
geometric hypotheses of \cref{lem:preliminary-short}, without its
upper width restriction. Put $q_\pm=B(\pm a/2)$,
$e_R=A(a)$, and $e_L=A(-a)$. Suppose
\begin{equation}\label{eq:preliminary-endpoint-laws}
 e_R=\tfrac12+\tfrac h4+\tfrac{3q_+-q_-}{4},\qquad
 e_L=\tfrac12-\tfrac{3h}{4}+\tfrac{3q_--q_+}{4}.
\end{equation}
Then the score is less than $41/50$.
\end{lemma}
\begin{proof}
Here use the original supports
$u=rh_U(\pi/4)-1$, $v=rh_U(3\pi/4)-1$,
$m=(u+v)/2$, $w_R=a-u$, and $w_L=a-v$.
The affine roof extension is $L_h(x)=1-h/2+(h/a)x$.
Concavity and the height bound imply
\[
 a/2\le u\le a,\quad a/2-h\le v\le a-3h/2,\quad
 u\ge a+e_R-1,\quad v\ge a+e_L-1.
\]
The majorant $\min(1,L_h,1+u-x,1+v+x)$ has exterior area
\[
 O_0=a-\frac58ah-\frac12
       \left[w_R^2+\frac{(w_L-3h/2)^2}{1-h/a}\right].
\]
The endpoint equations give
$e_R+e_L=1-h/2+(q_++q_-)/2\ge1-h/2$. Thus
\begin{equation}\label{eq:wide-endpoint-budget}
 w_R+w_L\le1+h/2,\qquad m\ge a-1/2-h/4.
\end{equation}
The diagonal tent has its apex in $J$, and its integral is
\[
 N_{45}=(m-k)^2-\tfrac12
       [(a/2-k-w_R)_+^2+(a/2-k-w_L)_+^2].
\]

Write $b_0=a/2-k$. For $t\ge b_0$ with $3/2-t/a\ge0$,
\[
 w^2-(b_0-w)_+^2\ge2tw-t^2,
\]
\[
 \frac{(w-3h/2)^2}{1-h/a}-(b_0-w)_+^2
 \ge2tw-t^2+(-3t+t^2/a)h.
\]
The second follows from the exact remainder
$[w-t-h(3/2-t/a)]^2/(1-h/a)$; if $w<b_0$, this remainder
is at least $(b_0-w)^2$. Hence
\begin{equation}\label{eq:wide-relaxation}
 P\le a-2t(a-m)+t^2-(m-k)^2
       -\left[\tfrac58a-\tfrac32t+\frac{t^2}{2a}\right]h.
\end{equation}
For $a\le1+2k$ take $t=1/2$. This yields
\[
 P\le m+1/4-(m-k)^2-D(a)h,\qquad
 D(a)=5a/8-3/4+1/(8a)\ge21/64.
\]
A hypothetical score at least $41/50$ forces
\begin{equation}\label{eq:wide-stability}
 a<9/5,\qquad h<1/4,\qquad a+h<37/20.
\end{equation}
Here are the scalar checks, retaining all width ranges.
For $9/5\le a\le1+2k$, substitute
$m\ge a-1/2-h/4$; the expression decreases in $h$, and its
value at $h=0$ decreases in $a$, giving
$P\le(23/5)k-57/50<41/50$.
For $a\ge1+2k$, take $t=b_0$ in
\cref{eq:wide-relaxation}; the same substitution gives
$P\le(5/2)k-1/4<41/50$.
For $a<9/5$ and $h\ge1/4$, the first relaxation gives
$P\le(163/40)k-2787/3200<41/50$.
Finally, if $a+h\ge37/20$ and $h<1/4$, substituting
$m\ge27/20-5h/4$ and completing the square gives
$P\le7529/25600+(101/80)k<41/50$.
Each strict comparison follows from $k<29/70$.

Put $z=h/a$ and $q=(v+h/2)/(1-z)$, so $a/2\le q\le a$.
The near-vertical deficits of the majorant are
\[
 \delta_L=c(1-h/2)+(s-cz)q-h_U(5\pi/8),\qquad
 \delta_R=c+su-h_U(3\pi/8).
\]
The middle points show $\delta_L\le(s-cz)(q-a/2)$ and
$\delta_R\le s(u-a/2)$, so cutting the two affine or flat
pieces removes whole charged triangles. Their total exterior
payment is at least $r(\delta_L^2+\delta_R^2)$.
The extra left deficit from tilt is
\[
 \tau=(c+sv-c(1-h/2)-(s-cz)q)/c
      =(1-k)h(1/2+q/a)\le\tfrac32(1-k)h.
\]

The $\pi/8$ tent has ascending wall
$b_R+kx$ and descending companion lower bound $B_R-x/k$, where
\[
 b_R=1-1/c+kv-\delta_L/c-\tau,\qquad
 B_R=e_R+a/k-1/s.
\]
In the reflected coordinate the corresponding $\pi/8$-type tent
at $3\pi/8$ has
$b_L=1-1/c+ku-\delta_R/c$ and $B_L=e_L+a/k-1/s$.
Their crossovers with the descending diagonal baselines are
\[
 x_L^R=(u-k-b_R)/r,\quad x_L^L=(v-k-b_L)/r,
\]
and their additional peaks are
\[
 Z_R=(B_R+b_R)/2-u+k,\qquad
 Z_L=(B_L+b_L)/2-v+k.
\]
The crossover offsets from the two diagonal apices are
$km+2s-1+(\delta_L/c+\tau)/r$ and
$km+2s-1+\delta_R/(rc)$, which are positive by
\cref{eq:wide-endpoint-budget,eq:wide-stability}.
The tentative peaks lie in $J$:
the middle support bounds reduce this to
$e_R+h+(3-r)a/2\le1+4s$ and
$e_L+(3-r)a/2\le1+4s$.
For the first use $a+h<37/20$, $h<1/4$, and
\[
 16s^2-(53/20-4r/5)^2=(-121+96r)/400>0;
\]
the second follows from $a<9/5$, $e_L\le1$, and $4s>3/2$.

Truncate each half-triangle at its diagonal baseline zero.
Set
\[
 \begin{gathered}
 F_R=r(u-k-x_L^R),\quad F_L=r(v-k-x_L^L),\\
 \widehat Z_Q=\min(Z_Q,F_Q),\quad
 \ell_Q=(Z_Q-F_Q)_+\quad(Q=R,L).
 \end{gathered}
\]
The two retained triangles are disjoint, lie in $J$, and add niche
area at least $[(\widehat Z_R)_+^2+(\widehat Z_L)_+^2]/(2r)$.
Their truncation lengths obey
\[
 \begin{split}
 2\ell_R&\le[e_R+2rw_R+h-(3k/2)a-1-4s+2rk]_+,\\
 2\ell_L&\le[e_L+2rw_L-(3k/2)a-1-4s+2rk]_+.
 \end{split}
\]
To bound these expressions we need no box estimate for a partial
niche. Solving \cref{eq:preliminary-endpoint-laws} using
$e_R\le1$, $e_L\le1-3h/2$, and $q_\pm\ge0$ gives directly
\[
 e_R\ge1/3+h/2+(2/3)q_+,\qquad
 e_L\ge1/3-2h/3+(2/3)q_-.
\]
In particular $w_R\le2/3-h/4$ and $w_L\le2/3+3h/4$.
Using also $e_Q\le1-w_Q$, $a\ge8/5$, and
\[
 (2r-1)\tfrac23-\tfrac{12}{5}k-4s+2rk<0
\]
in the displayed truncation bounds gives
\[
 \ell_R\le(5-2r)h/8,\quad
 \ell_L\le3(2r-1)h/8,\quad \ell_R+\ell_L\le h.
\]
The constant's sign follows from $r<10/7$, $k<3/7$, $s>3/8$.

Adding the two peak formulas, using $e_R+e_L\ge1-h/2$,
$\tau\le3(1-k)h/2$, and the last truncation bound, gives
\[
 \widehat Z_R+\widehat Z_L
 \ge Q(a,m,h)-(\delta_L+\delta_R)/(2c),\qquad
 Q=\tfrac{12}{5}a-\tfrac85m-\tfrac{11}{8}-\tfrac{17}{10}h.
\]
If $D_0=r(\delta_L^2+\delta_R^2)$ and
$G_0=[(\widehat Z_R)_+^2+(\widehat Z_L)_+^2]/(2r)$,
Cauchy--Schwarz gives $Q_+^2\le8(D_0+G_0)$.
The geometric payments therefore yield
\[
 P\le\widetilde G(a,m,h)
 =m+1/4-(m-k)^2-\tfrac{21}{64}h-\tfrac18Q_+^2.
\]
This is jointly concave. At $(a_0,m_0,h_0)=(8/5,11/10,0)$,
its gradient is
\[
 G_a=-423/1000,\qquad G_m=2k-459/500,\qquad G_h=-57/2000.
\]
By \cref{eq:wide-endpoint-budget}, a feasible displacement is
\[
 (a-a_0)(1,1,0)
 +(m-a+1/2+h/4)(0,1,0)+h(0,-1/4,1).
\]
All coefficients are nonnegative. All three directional derivatives
are negative; the last is $201/1000-k/2<0$, since $k>41/100$.
The supporting plane gives
\[
 P\le\widetilde G(a_0,m_0,0)
   =\frac{21}{5}k-\frac{295081}{320000}
   <\frac{261719}{320000}<\frac{41}{50}.
\]
This contradiction proves the lemma.
\end{proof}

\subsection{The intermediate tilt bound}

\begin{lemma}[One-angle tilt estimate]\label{lem:preliminary-tilt}
For $1\le a\le8/5$, $0\le h<1/2$, and the same middle roof,
a barrier dominating the actual diagonal tent gives
\begin{equation}\label{eq:intermediate-tilt}
 P\le a-\tfrac58ah-\frac{2(a-3h/4-k)^2}{4-h/a}
  =1-\frac{(a-r)^2}{2}
       -\frac{2a}{4a-h}(k+h/2)^2.
\end{equation}
A score at least $41/50$ consequently forces $h<17/50$.
\end{lemma}
\begin{proof}
Use the lifted cap and its diagonal supports $u,v$ from the short
proof. Concavity of the original roof imposes $v\le a-h/2$.
In addition to the horizontal majorant, the lifted left wing lies
below $1+h/2+(h/a)x$. Its cut removes exterior area
\[
 D_{\rm tilt}=\frac{h(v-a/2)^2}{2(a-h)}.
\]
Indeed the first triangle has area
$(h/a)(v-a/2)^2/2$, and the second has this value times
$(h/a)/(1-h/a)$; both are charged because $v\le a-h/2$.
The score is therefore bounded by
$F_h(a,u,v)-D_{\rm tilt}$, a concave function on the entire
square $[a/2,a]^2$.

Put
\[
 D_*=\frac{2(a-3h/4-k)}{4-h/a},\quad
 u_*=a-D_*,\quad v_*=a-h/2-(1-h/a)D_*.
\]
This is its unclipped stationary point. Its two coordinates are
strictly inside the square, since $0<D_*<a/2$.
Its diagonal apex has height $D_*>0$. The right clip vanishes since
\[
D_*-(a/2-k)=\frac{2k-h(1+k/a)}{4-h/a}>0,
\]
using $h<1/2$, $a\ge1$, and $r>4/3$.
The left clip is smaller still. Substituting this global critical
point proves the first formula in \cref{eq:intermediate-tilt}.
The second is its expansion using $k^2=1-2k$.

The bound decreases with $h$. Dropping the nonpositive width square
and using $a\le8/5$ gives
\[
 P\le1-\frac{16/5}{32/5-h}(k+h/2)^2.
\]
For $h\ge17/50$, the right side is at most its value at $17/50$.
Since $k>41/99$, this is strictly below
\[
 1-\frac{16/5}{32/5-17/50}(41/99+17/100)^2
 =\frac{304326697}{371212875}<\frac{41}{50}.
\]
The final rational gap is $135721/742425750>0$.
\end{proof}

\begin{corollary}[Canonical range]\label{lem:canonical-range}\label{lem:preliminary-cuts}
Every canonical full maximum satisfies
\[
 \frac{1001}{2000}<C<\frac45,\qquad
 |f_U(C)-f_U(-C)|<\frac{17}{50},
 \quad I_U=[-2C,2C].
\]
The same conclusions hold for every canonical proper joint maximum
whose angle exceeds $\arctan(8/3)$.
\end{corollary}
\begin{proof}
The reference competitor and \cref{lem:constant-bounds} give
score $M/2>41/50$. Width coercivity gives $a=2C<3$.
The positive-pressure and top-localization results give height one,
an affine middle roof, and lower middle height greater than $1/2$.
Orient its higher middle endpoint to the right. In the proper case
reflect the actual barrier as well. It remains nonnegative and
dominates all three tents, since
$\arctan(8/3)>3\pi/8$. The endpoint equations are precisely
\cref{eq:preliminary-endpoint-laws} after relabeling.
Apply \cref{lem:preliminary-short,lem:preliminary-wide} to obtain
$1001/1000<a<8/5$, then \cref{lem:preliminary-tilt} to obtain
the height-drop bound.
\end{proof}
